\documentclass[11pt,a4paper]{article}
\usepackage[T1]{fontenc}
\usepackage[utf8]{inputenc}
\usepackage[spanish]{babel}
\usepackage{amsmath,amssymb}
\usepackage[margin=2.5cm]{geometry}
\usepackage[hidelinks]{hyperref}
\newcommand{\B}{\{0,1\}}
\title{Murdoku: modelo del caso de uso\\Verificar un puzzle}
\author{Documentaci\'on de la implementaci\'on}
\date{}

\begin{document}
\maketitle

\section{Alcance y fuentes}

El caso de uso \texttt{verify-puzzle}, ejecutado con
\texttt{murdoku-checker}, recibe mapa, cat\'alogo de elementos, personas
y pistas. Busca hasta dos soluciones para distinguir imposibilidad,
unicidad y ambig\"uedad. El mapa y los muebles son \textbf{datos fijos}.
No se genera una nueva distribuci\'on ni se optimiza la dificultad.

Las fuentes principales son:
\begin{itemize}
  \item \path{murdoku_checker\apps\solution_checker\variable_builder.py}.
  \item \path{murdoku_checker\apps\solution_checker\constraint_builder.py}.
  \item \path{murdoku_checker\apps\solution_checker\model.py}.
  \item \path{murdoku_checker\apps\solution_checker\solution_checker.py}.
  \item Las estrategias de \path{murdoku_checker\domain\rules}.
\end{itemize}
Las f\'ormulas siguientes describen las restricciones l\'ogicas;
CP-SAT puede emplear literales auxiliares equivalentes para expresarlas.

\section{Datos y variables}

Sean $P$ las personas, $V$ las celdas existentes, $F$ los tipos de
elemento y $\mathcal R$ las pistas recibidas. Para $c\in V$ son datos:
su fila $i(c)$, columna $j(c)$, habitaci\'on $\kappa(c)$ y conjunto
de elementos $F(c)$. Sea $o_f=1$ si el elemento $f$ es ocupable.

El conjunto de celdas disponibles es
\[
 U=\{c\in V:\ o_f=1\text{ para todo }f\in F(c)\}.
\]
Una celda sin elementos pertenece a $U$. Si contiene varios, todos
deben ser ocupables. Las referencias a personas, habitaciones y
elementos se validan antes de resolver; una referencia inexistente
es un error de entrada.

Se construye una variable por persona y celda disponible:
\[
 q_{pc}\in\B\qquad(p\in P,\ c\in U),
\]
donde $q_{pc}=1$ indica que $p$ est\'a en $c$. No se crean variables
de colocaci\'on para celdas bloqueadas.

\section{Modelo base de colocaci\'on}

Las restricciones comunes a todas las pistas son
\begin{align}
 \sum_{c\in U}q_{pc}&=1 &&\forall p\in P,\\
 \sum_{\substack{p\in P,\ c\in U\\i(c)=i}}q_{pc}&\leq 1
      &&\forall\text{ fila }i,\\
 \sum_{\substack{p\in P,\ c\in U\\j(c)=j}}q_{pc}&\leq 1
      &&\forall\text{ columna }j.
\end{align}
Si una persona carece de celdas disponibles, su primera igualdad
es $0=1$, de modo que el modelo es infactible.

Se definen las expresiones, no nuevas decisiones independientes,
\[
 Y_p=\sum_{c\in U}i(c)q_{pc},\qquad
 X_p=\sum_{c\in U}j(c)q_{pc},\qquad
 K_p=\sum_{c\in U}\kappa(c)q_{pc}.
\]
Para $K_p$ se asigna a cada nombre de habitaci\'on un entero distinto.
La colocaci\'on exactamente una vez hace que las expresiones sean
las coordenadas y la habitaci\'on elegidas.

\paragraph{Diferencia respecto al generador.}
El checker no vuelve a construir habitaciones ni muebles y no impone
conectividad, tama\~nos de habitaciones o porcentajes de cobertura.
Tampoco se aplica el l\'imite del generador de como m\'aximo una
habitaci\'on contenida en una sola fila o columna: un mapa fijo
existente no se rechaza por contener varias habitaciones estrechas.
Tampoco se exigen rectangularidad ni proporci\'on m\'axima de 1.5;
estas son restricciones del generador, no del checker de mapas fijos.
Tampoco impone autom\'aticamente que la v\'ictima comparta habitaci\'on
con exactamente otra persona. Esa condici\'on pertenece al generador
de soluciones. La marca de v\'ictima y los dem\'as atributos personales
no a\~naden restricciones por s\'i solos. Tampoco se exige cobertura
narrativa de las personas: es una restricci\'on del maestro generador.

\section{Modelo de las pistas}

Cada pista $r$ tiene un \'unico protagonista $p(r)$, el sujeto de su
cl\'ausula. Una persona mencionada como referencia es un objetivo,
no un protagonista adicional. Esto no impide que la pista restrinja
tambi\'en la posici\'on de esa persona.

Sea $\varphi_r(q)$ el predicado de verdad de la pista. Aplicar una
pista al modelo equivale a imponer $\varphi_r(q)=1$. Cuando se
compone dentro de otra, su literal $b_r\in\B$ cumple
\[
 b_r=1\Longleftrightarrow\varphi_r(q).
\]
La equivalencia se implementa en ambos sentidos, no como una sola
implicaci\'on. Las conjunciones y disyunciones vac\'ias tienen los
valores l\'ogicos usuales: verdadero y falso, respectivamente.

\subsection{Pertenencia: \texttt{in}}

Para un objetivo habitaci\'on $k$, sea
$A_r=\{c\in V:\kappa(c)=k\}$; para un objetivo elemento $f$, sea
$A_r=\{c\in V:f\in F(c)\}$. La pista exige
\[
 \sum_{c\in U\cap A_r}q_{p(r)c}=1.
\]
Con \texttt{is\_unique=True} se exige adem\'as
\[
 q_{uc}=0\qquad
 \forall u\in P\setminus\{p(r)\},\ c\in U\cap A_r.
\]
La exclusividad se aplica a todas las celdas de la habitaci\'on o
de ese tipo de elemento, no solo a un ejemplar concreto.

\subsection{Coordenada fija: \texttt{in\_coordinate}}

Para un valor $a$ y eje fila o columna, respectivamente:
\[
 Y_{p(r)}=a\qquad\text{o}\qquad X_{p(r)}=a.
\]
Los datos usan coordenadas desde cero; la descripci\'on textual de la
pista muestra la coordenada m\'as uno.

\subsection{Adyacencia: \texttt{next\_to} y \texttt{not\_next\_to}}

Para un tipo de elemento $f$, se precalcula
\[
 A_f=\{c\in V:\ \exists u\in V,\ f\in F(u),\
 |i(c)-i(u)|+|j(c)-j(u)|=1,\ \kappa(c)=\kappa(u)\}.
\]
\texttt{next\_to} exige que la persona est\'e en $A_f$;
\texttt{not\_next\_to}, que est\'e en $V\setminus A_f$.
No se admiten diagonales ni adyacencia a trav\'es de una pared.
La segunda pista es la negaci\'on exacta de la primera, aunque no
haya ninguna celda con el elemento.

\subsection{Relaci\'on de habitaci\'on: \texttt{room\_relation}}

Dados el sujeto $p$ y la persona objetivo $u$:
\[
 \varphi_r(q)=
 \begin{cases}
  [K_p=K_u],&\texttt{same},\\
  [K_p\neq K_u],&\texttt{different}.
 \end{cases}
\]
Los corchetes indican el valor booleano del predicado.

\subsection{Posici\'on relativa: \texttt{relative\_position}}

Sea $C_p=Y_p$ para norte o sur, y $C_p=X_p$ para este u oeste.
El signo es $\sigma=1$ para sur o este y $\sigma=-1$ para norte
u oeste. Para una diferencia $\delta$:
\[
 M_r(\delta)=
 \begin{cases}
  [\sigma\delta>0],&\text{sin distancia exacta},\\
  [\delta=\sigma h],&\text{con distancia exacta }h\geq1.
 \end{cases}
\]
Con objetivo persona $u$, $\varphi_r(q)=M_r(C_p-C_u)$.
Solo se compara el eje indicado: no se exige coincidencia en el
otro eje ni pertenencia a la misma habitaci\'on.

Con objetivo elemento $f$, la referencia es \emph{otra persona
existente}, no la coordenada de un mueble vac\'io. Sea $A$ el conjunto
de celdas con $f$ si \texttt{feature\_relation=in}, o $A_f$ de la
secci\'on anterior si vale \texttt{next\_to}. Entonces:
\[
 \varphi_r(q)=
 \bigvee_{\substack{u\in P\setminus\{p\},\ c\in U,\ d\in U\cap A\\
 M_r(C(c)-C(d))=1}}
       (q_{pc}\land q_{ud}).
\]
$C(c)$ es la coordenada del eje correspondiente de $c$.
La implementaci\'on crea literales de conjunci\'on para los pares
compatibles y una disyunci\'on sobre ellos.

\subsection{Disyunci\'on: \texttt{any\_of}}

Para las condiciones $r_1,\ldots,r_\ell$, con $\ell\geq2$:
\[
 \varphi_r(q)=\bigvee_{a=1}^{\ell}\varphi_{r_a}(q).
\]
Es una disyunci\'on inclusiva: varias condiciones pueden ser ciertas.
Se admiten condiciones anidadas al verificar y todas deben tener el
mismo sujeto que la regla exterior. El checker no convierte esta
expresi\'on en una disyunci\'on exclusiva.

\section{Modelo completo y enumeraci\'on}

El conjunto de soluciones es
\[
 \mathcal F(\mathcal R)=
 \{q:\ q\text{ satisface el modelo base},\
       \varphi_r(q)=1\ \forall r\in\mathcal R\}.
\]
No hay objetivo. \texttt{check\_solutions} activa
\texttt{enumerate\_all\_solutions} y detiene el callback cuando encuentra
dos soluciones. Conserva solo la primera colocaci\'on:
\begin{center}
\begin{tabular}{cl}
\hline
Recuento & Interpretaci\'on\\
\hline
$0$ & No hay soluci\'on.\\
$1$ & Soluci\'on \'unica, tras agotar la b\'usqueda.\\
$2$ & Al menos dos soluciones; no es el total exacto.\\
\hline
\end{tabular}
\end{center}
Un estado no concluyente del solver produce \texttt{RuntimeError}.
Un l\'imite de enumeraci\'on igual a uno no basta para demostrar
unicidad; el comando usa siempre dos.

\section{Variante usada por el generador: excluir un objetivo conocido}

La funci\'on \texttt{find\_alternative} recibe adem\'as una colocaci\'on
$S^\star$, con celda $c_p^\star$ para cada persona. Primero comprueba
que usa exactamente las personas y el mapa de entrada, que no repite
filas ni columnas, que ocupa celdas disponibles y que satisface las
pistas. Despu\'es a\~nade:
\[
 \bigvee_{p\in P}\neg q_{p,c_p^\star},
 \qquad\text{equivalentemente}\qquad
 \sum_{p\in P}q_{p,c_p^\star}\leq |P|-1.
\]
La cl\'ausula exige cambiar al menos una persona, no todas. Se busca
una \'unica alternativa:
\begin{itemize}
  \item Si aparece, devuelve $S_{\mathrm{alt}}\neq S^\star$.
  \item Si no existe ninguna, devuelve \texttt{None}; como $S^\star$
        es factible, queda demostrada su unicidad.
  \item Si el objetivo es inv\'alido o el estado no es concluyente,
        se comunica un error, no una prueba de unicidad.
\end{itemize}
Esta variante reutiliza el mismo modelo base y las mismas estrategias
de pistas. No modifica habitaciones, muebles ni personas.

\end{document}
